% =============================================================================
%  slushpile cover letter template
%
%  Build:  latexmk -xelatex cover_letter.tex && latexmk -c
%  Needs:  XeLaTeX (for fontspec). Same font set and fallback as resume.tex.
%          The house fonts are vendored: `python3 scripts/install_fonts.py`
%          installs them. Skipping it falls back to DejaVu and still builds.
%          The FONTS block in templates/resume.tex explains why it works this
%          way rather than by path.
%
%  This is the visual companion to resume.tex: same palette, same fonts, same
%  header treatment. Send them together. A letter in a different typeface than
%  the resume reads as two documents from two sources, and the pair is the only
%  thing in the application a screener sees side by side.
%
%  ---------------------------------------------------------------------------
%  WHAT THIS DOCUMENT IS FOR
%  ---------------------------------------------------------------------------
%  The resume answers "can this person do the job." The letter answers "why this
%  person, at this company, now." If a paragraph here could be moved into the
%  resume without loss, it belongs in the resume.
%
%  The body is written by the voice agent named in preferences.yaml under
%  `voice.agent`, not drafted here. This file is the container and the structure.
%  Check `voice.is_mine` before you send anything: if it is false, the letter is
%  in the shipped example voice rather than yours.
%
%  ---------------------------------------------------------------------------
%  STRUCTURE, AND WHY IT IS THIS ORDER
%  ---------------------------------------------------------------------------
%   1. The hook. One or two sentences. Never background. "I have spent twelve
%      years doing X" is the most common opening in the bottom half of the pile
%      and a screener has read it forty times before reaching yours.
%   2. The thesis, made company-specific. Test it by swapping the company name
%      for a direct competitor. If the thesis survives the swap it is not
%      specific enough, and the hiring manager runs that same test by instinct.
%   3. Evidence, two to four paragraphs. One story with room to breathe beats
%      three summarized ones. Lead each paragraph with the claim, then the
%      mechanism, then the number.
%   4. The close. The thing you would build here first, or the habit you would
%      bring. Concrete. Not "I would welcome the opportunity to discuss."
%
%  ---------------------------------------------------------------------------
%  LENGTH
%  ---------------------------------------------------------------------------
%  One page is the convention and it is a real constraint: a second page of a
%  cover letter is read at a much lower rate than the first. Going over is a
%  decision to make deliberately, per application, not a thing to discover after
%  the build. Check the page count before you send.
%
%  ---------------------------------------------------------------------------
%  HOUSE STYLE
%  ---------------------------------------------------------------------------
%  No em dashes: comma, colon, semicolon, or parentheses instead. En dashes for
%  ranges. Contractions are fine and usually better -- this is the one document
%  in the application where a human voice is the point.
%
%  Placeholders are UPPERCASE so a forgotten one is visible on the page.
% =============================================================================

\documentclass[10pt,letterpaper]{article}

\usepackage[top=0.45in,bottom=0.4in,left=0.65in,right=0.65in]{geometry}
\usepackage{fontspec}
\usepackage{xcolor}
\usepackage{textcomp}
\usepackage{microtype}
\usepackage{hyperref}

% ---- palette: identical to resume.tex. Change both or neither. -------------
\definecolor{base900}{HTML}{1A1B1E}
\definecolor{ink}{HTML}{181B20}
\definecolor{muted}{HTML}{43505D}
\definecolor{faint}{HTML}{6B7480}
\definecolor{accent}{HTML}{B05B33}
\definecolor{accentink}{HTML}{7C3E20}

\hypersetup{
    colorlinks=true, urlcolor=accentink, linkcolor=accentink,
    pdftitle={NAME -- Cover Letter, ROLE}, pdfauthor={NAME}
}

% ---- fonts, with the same fallback so this builds on a bare TeX Live -------
\IfFontExistsTF{Public Sans}{%
    \setmainfont{Public Sans}[
      UprightFont={* Regular}, BoldFont={* Bold},
      ItalicFont={* Italic}, BoldItalicFont={* Bold Italic}]
    \newfontfamily\heavy{Public Sans}[UprightFont={* ExtraBold}]
}{%
    \setmainfont{DejaVu Sans}
    \newfontfamily\heavy{DejaVu Sans}[UprightFont={* Bold}]
}

\IfFontExistsTF{IBM Plex Mono}{%
    \setmonofont{IBM Plex Mono}[Scale=0.84, UprightFont={* Regular}, BoldFont={* SemiBold}]
}{%
    \setmonofont{DejaVu Sans Mono}[Scale=0.78]
}

\color{ink}
\linespread{1.04}
\setlength{\parindent}{0pt}
\setlength{\parskip}{5.5pt}
\emergencystretch=3em
\hyphenpenalty=400

\newcommand{\eyebrow}[1]{{\ttfamily\scriptsize\color{faint}\addfontfeatures{LetterSpace=4.0}\MakeUppercase{#1}}}

\pagestyle{empty}

\begin{document}

% ============================== HEADER ==============================
% Same contact line as the resume, character for character. A phone number that
% differs between the two documents is the kind of small inconsistency a
% fatigued reader notices and then starts looking for more of.
%
% The Re: line carries the exact role title and the requisition ID from the
% posting. Large employers route on the requisition, and a letter that names it
% lands in the right pile without a human deciding where it goes.

{\heavy\fontsize{19}{21}\selectfont\color{base900}NAME}\\[2pt]
{\ttfamily\footnotesize\color{muted}CITY, STATE \ \textbullet\ PHONE \ \textbullet\ EMAIL \ \textbullet\ \href{https://github.com/USER}{github.com/USER} \ \textbullet\ \href{https://linkedin.com/in/USER}{linkedin.com/in/USER}}\\[6pt]
\eyebrow{Re: ROLE TITLE \ \textbullet\ \ req REQUISITION ID (LOCATION) \ \textbullet\ \ DD Month YYYY}

\vspace{4pt}
{\color{accent}\hrule height 1.4pt}
\vspace{4pt}

% ============================== SALUTATION ==============================
% Name the hiring manager when you know it and can verify it. Do not guess: a
% letter addressed to the wrong person is worse than one addressed to nobody.
% "Hi," is the safe default and it costs you nothing. "To Whom It May Concern"
% costs you something.

Hi,

% ============================== HOOK ==============================
% The first line presents the angle, not the background. It lands when it makes
% a reader curious about someone they would not otherwise interview: a concrete
% detail, a reframe, or a provocation. Not a credential.
%
% Then, in the same paragraph, turn it toward this company specifically. Name
% something true about how they operate that you could not say about a
% competitor, and connect it to the problem in the posting.

THE HOOK. One or two sentences that reframe the problem in this role, stated as
a fact rather than a claim about yourself. THE TURN: the thing this company does
that a competitor does not, and why it makes the reframe theirs rather than
generic. THE CONNECTION back to the seat being filled.

% ============================== EVIDENCE ==============================
% Two to four paragraphs. Each opens with the claim, then gives the mechanism,
% then the number. A paragraph that opens with the number and works backward
% reads as a metrics dump.
%
% Every factual claim here must already exist in profile.md, and every number
% must match the "Numbers You Can Defend" table. A figure that appears as 2,400
% in the letter and 2,707 in the resume reads as carelessness, and a reader who
% catches one inconsistency starts hunting for others.
%
% Link the public artifacts inline. A stranger who clicks one and finds it real
% has verified you in thirty seconds without calling a reference.

Here is what that looks like in my own work. \href{https://github.com/USER/PROJECT}{PROJECT}
is THE ONE-LINE DESCRIPTION. THE DECISION YOU MADE, and the failure it prevents.
The version that would have been easier is THE EASIER WRONG THING, and it would
have been WHY IT WAS WRONG.

THE MAIN STORY, from stories.md. Give it the space: scene, then the
complication, then what you personally did as distinct from what the team did,
then the resolution with its number. This is the paragraph a reader remembers,
and the reason it works is the moment in it, not the summary of it. Check the
"Used in" field before reusing a story on a second application to this company.

THE PARAGRAPH THAT ANSWERS THE ROLE'S HARDEST REQUIREMENT. If the role analysis
found a gap that the posting weights heavily, this is where the nearest real
thing you have done goes. Do not name the gap and do not apologize for it. The
pipeline does not volunteer weaknesses to employers: recorded gaps exist so the
fit scoring is honest, not so the letter can confess them.

% ============================== CLOSE ==============================
% One paragraph. The specific thing you would do here, or the habit you would
% bring, named concretely enough that it could be wrong. A close that could
% appear at the bottom of any letter is a wasted paragraph in the position a
% reader is most likely to actually finish.

THE CLOSE. The one practice you would bring, described so specifically that
someone could disagree with it. Then why it matters here rather than in general.

\vspace{2pt}
NAME\\[1pt]
\eyebrow{WORK AUTHORIZATION \ \textbullet\ \ CERTIFICATIONS}

\end{document}
